\documentclass[11pt]{article}

\usepackage[margin=1in]{geometry}
\usepackage{amsmath}
\usepackage{amssymb}

\usepackage{graphicx}
\usepackage{bm}
\usepackage{booktabs}
\usepackage{array}
\usepackage{amsthm}
\usepackage{enumerate}
\usepackage[numbers,sort&compress]{natbib}
\usepackage{hyperref}
\usepackage{lineno}

\hypersetup{
  colorlinks=true,
  linkcolor=blue,
  citecolor=blue,
  urlcolor=blue
}

\begin{document}

\title{A Calibrated Vortex-Ring Toy Model for the Fine-Structure Ratio}

\author{Steve Bico Mujjabi, MD\thanks{msbico@gmail.com}\\
VuraLabs, Kampala, Uganda\\
Independent researcher\\
ORCID: \href{https://orcid.org/0009-0001-0556-5516}{0009-0001-0556-5516}\\
\url{https://github.com/bampita-bico/CDFD-Part-I-Release}}

\date{July 2026}

\maketitle

\begin{abstract}
The electromagnetic fine-structure constant is a measured input to the Standard
Model. This paper records a conditional vortex-ring toy model in which a chosen
dimensionless stiffness is set so that a specified energy functional has its
minimum at the known inverse fine-structure ratio. The numerical recovery is
therefore a calibration check: the reported residual is the root solver's
convergence tolerance, not a prediction error. The accompanying length-scale
calculation gives $R/a=\hbar c/e^2=1/\alpha$ by the definitions of the reduced
Compton wavelength, classical electron radius, and $\alpha$; it is not an
independent confirmation.

The transport-medium, capacity, and memory language remains a hypothesis. The
calculus of the stated energy functional is valid conditional on its assumptions,
but the stiffness has no independent physical derivation and the model provides
no established vacuum mechanism or fine-structure prediction. A correction
status notice in the Part I release gives the binding interpretation boundary.
\end{abstract}

\noindent\textbf{Keywords:} Constraint-Driven Field Theory, Constraint-Driven Flux Dynamics, adaptive operating ratio, vacuum structure, Mujjabi tests, reproducible physics, fine-structure constant, vortex stability

\vspace{0.75em}
\linenumbers

\section*{Project Context}
This manuscript constitutes Part I (Paper 01) of the multi-part series
\textit{Constraint-Driven Flux Dynamics (CDFD)}. It is the opening mathematical
and physical paper in the Part I stack. Downstream Papers 02--12 are organized in the same article-class release family and cite this manuscript by title until paper-specific DOIs are assigned.

\section*{Symbol Conventions}

\noindent\textbf{CDFD physical-vacuum symbol convention.}
\begin{center}
{\small
\begin{tabular}{p{0.16\linewidth}p{0.33\linewidth}p{0.41\linewidth}}
\hline
Symbol & Part I meaning & Physical reading \\
\hline
$J$ or $\Phi$ & Driving flux or throughput & Transport current, field throughput, circulation drive, or generic flux; a subscript fixes the exact channel \\
$C$ & Active constraint or capacity burden & Vacuum transport capacity, geometric confinement, stiffness, boundary resistance, or load-bearing limit \\
$S$ & Surface responsiveness & Local ability of the vacuum or modeled surface to reconfigure under drive \\
$M_s$ & Structural memory & Retained geometry, phase history, stiffness history, or accumulated constraint state \\
$\Psi_s$ & Adaptive operating ratio $(J/C)\cdot S\cdot M_s$ or $(\Phi/C)\cdot S\cdot M_s$ & Local bookkeeping gauge for constrained operation; it is not a new universal constant \\
$R,a,\chi$ & Ring radius, core radius, and aspect ratio $R/a$ & Geometry used to connect vortex structure to dimensionless electromagnetic scales \\
$\rho_0,c,\kappa$ & Vacuum density scale, propagation speed, and transport stiffness & Medium parameters used in stability, pressure, and self-consistency calculations \\
\hline
\end{tabular}
}
\end{center}

\noindent\textbf{Greek-symbol convention.}
\begin{center}
{\small
\begin{tabular}{p{0.16\linewidth}p{0.33\linewidth}p{0.41\linewidth}}
\hline
Symbol & Local role & Interpretation rule \\
\hline
$\alpha$ & Fine-structure constant unless a local equation states otherwise & In the Part I physics papers, un-subscripted $\alpha$ means the electromagnetic fine-structure constant \\
$\beta,\gamma$ & Local response, stiffness, coupling, or correction coefficients & Equation-local coefficients; subscripts bind them to the named mode, channel, or approximation \\
$\theta,\phi,\Phi_c$ & Phase, orientation, or chirality angles & Used for torus-knot orientation, vacuum phase, or chiral phase geometry as defined in the nearby equation \\
$\delta,\epsilon$ & Perturbation, residual, tolerance, or small-offset scales & Local stability margins, numerical residuals, or experimental tolerances \\
$\lambda,\mu,\nu$ & Rate, mass, mode, or frequency labels & Meaning is fixed by the equation, subscript, and physics context where the symbol appears \\
$\rho,\sigma,\tau$ & Density, spread/noise, or time-scale terms & Local density, variance, stress, relaxation, or transport-memory quantities \\
$\Omega,\Lambda$ & State/domain space or integrated measure where defined & Reserved for explicitly defined state spaces, spectra, or synthesis-level quantities \\
\hline
\end{tabular}
}
\end{center}


\section{Introduction}

Richard Feynman once wrote that the fine-structure constant ``has been a mystery
ever since it was discovered more than fifty years ago, and all good theoretical
physicists put this number up on their wall and worry about it''~\cite{Feynman1964}.
That was in 1964. Sixty years later, the mystery remains open. The
Standard Model of particle physics accepts $\alpha$ as an experimental input --- the
coupling constant of quantum electrodynamics~\cite{Schwinger1948}. It
does not explain it.

The number itself is $\alpha = e^2/\hbar c \approx 1/137.036$. It controls the
strength of the electromagnetic interaction. If it were much larger, atoms would
not hold together. If it were much smaller, stars would not ignite. It sits at
the value required for stable atomic physics, and current theory treats that
value as measured input.

Previous attempts to derive $\alpha$ have not produced a broadly accepted
mechanism. Eddington predicted
$1/\alpha = 136$ from quantum combinatorics~\cite{Eddington1929} --- wrong value,
no physical mechanism. Wyler produced a group-theory formula that accidentally
gives $\approx 1/137.03$~\cite{Wyler1969}, but it provides no dynamical explanation
for why the vacuum selects that geometry. More recent stochastic approaches remain
incomplete~\cite{BeckMackey2004}. The proposal below is a toy construction: it
fixes an energy functional and examines the algebraic consequences of
calibrating one of its parameters to a known constant.

Within that construction, the vacuum is treated as a hypothetical transport
medium with finite capacity. A vortex ring is a modelling ansatz, not an
established physical constituent of the vacuum. The classical theory of vortex
motion was founded by Helmholtz~\cite{Helmholtz1858} and developed by
Kelvin~\cite{Thomson1867}, Lamb~\cite{Lamb1932}, and Saffman~\cite{Saffman1992}.
The calibration below chooses $\kappa$ so that its equilibrium condition is met
at $\chi=1/\alpha$. It does not derive $\alpha$, determine a vacuum stiffness,
or establish the physical interpretation of the ansatz.

\subsection{Headline numerical result}

The public calculation takes the CODATA value
\[
  \alpha^{-1}_{2022}=137.035999177(21),
\]
as its target. Solving Eq.~\eqref{eq:kappa_for_chi} for the stiffness that gives
that target, then solving the same equilibrium condition with the fitted
stiffness, returns
\[
  \chi_\mathrm{recovered}=137.035999177,
  \qquad
  \alpha_\mathrm{recovered}=0.007297352566.
\]
The reported $1.19\times10^{-12}$ relative difference is the bisection
convergence tolerance. The length-scale calculation is an algebraic identity,
not a separate check, because the definitions already give
$R/a=\hbar c/e^2=1/\alpha$. Neither calculation validates the vacuum-medium
hypothesis.

\subsection{Public runtime and paper-local reproducibility}

The CDFD Runtime is now public as the reusable runtime for the
flow--constraint--memory state grammar used across the broader project. This
paper does not depend on hidden runtime behavior. Its numerical claims are
reproduced by the paper-local supplementary script and notebook in this Part I
archive. The runtime provides the general implementation path for future
domain-level simulations; the Part I scripts remain the audited source for the
figures, tables, and checks cited here.

\subsection{Named hypotheses and notation}

For clarity, the candidate constructs are named here at the point where their
symbols are first introduced. These labels do not confer the status of physical
laws, and they must not be used as evidence for downstream claims.

\begin{center}
\small
\begin{tabular}{p{0.25\textwidth}p{0.38\textwidth}p{0.24\textwidth}}
\toprule
\textbf{Name} & \textbf{Statement inside CDFT} & \textbf{Status}\\
\midrule
Mujjabi Capacity Hypothesis & Local transport becomes nonlinear as $J/C\rightarrow1$; vortex states regulate overload. & Unvalidated hypothesis\\
Mujjabi Adaptive Operating Ratio & $\Psi_s=(J/C)\,S\,M_s$ tracks flow, capacity, responsiveness, and memory. & Core notation\\
Calibrated Stability Minimum & The selected energy functional has a minimum at the target $\chi=R/a=1/\alpha$ after fitting $\kappa$. & Conditional calculus\\
Vacuum Memory Hypothesis & $M_s$ is a local constraint-history state with finite recovery time $\tau_M$. & Unvalidated hypothesis\\
\bottomrule
\end{tabular}
\end{center}

These hypotheses create failure routes. If $J/C$ has no predictive value, the
capacity hypothesis fails. If $M_s$ leaves no measurable residual in controlled
high-flux tests, the memory hypothesis fails. If core-profile changes destroy a
calibrated fixed point without a principled, independently determined
correction, the model adds no explanatory value.

\subsection{Symbol ledger}

The main symbols used in Paper~I are:
\begin{center}
\small
\begin{tabular}{p{0.18\textwidth}p{0.68\textwidth}}
\toprule
\textbf{Symbol} & \textbf{Meaning in this manuscript}\\
\midrule
$\alpha$ & electromagnetic fine-structure constant, written here as $e^2/(\hbar c)$ in the manuscript's electromagnetic-unit convention; $\alpha^{-1}$ is the CODATA inverse fine-structure ratio\\
$R$ & vortex-ring radius, identified here with the reduced Compton scale $\hbar/(mc)$\\
$a$ & vortex-core radius, identified here with the classical electron radius $e^2/(mc^2)$ in the manuscript's electromagnetic-unit convention\\
$\chi$ & dimensionless vortex aspect ratio $R/a$\\
$\Gamma$ & circulation of the vortex state\\
$\beta$ & dimensionless vortex-core profile constant; $\beta=7/4$ for the Lamb solid-rotation baseline\\
$\kappa$ & dimensionless transport stiffness in the normalized vortex energy\\
$J$ & local flux magnitude through the vacuum medium\\
$J_\mathrm{crit}$ & critical flux capacity of the medium\\
$S$ & local responsiveness or adaptive surface factor\\
$M_s$ & retained structural-memory factor of the medium\\
$\Psi_s$ & adaptive operating ratio $(J/J_\mathrm{crit})SM_s$ in Paper~I\\
$\tau_M$ & candidate relaxation time for vacuum-memory decay\\
$\mu$ & candidate coupling strength between overload and memory accumulation\\
\bottomrule
\end{tabular}
\end{center}


\section{The Adaptive Vacuum Medium}

\subsection{Governing equation}

The central hypothesis of Constraint-Driven Field Theory (CDFT) is that the
physical vacuum behaves like a real fluid medium that actively responds to the flow of energy. The idea that the vacuum is a
physical medium with definite properties is not new --- Dirac himself argued for
it in 1951~\cite{Dirac1951}, and modern treatments of vacuum condensates pursue
related ideas~\cite{Consoli2008}. What CDFT adds is a transport capacity constraint,
memory-bearing state variables that track the medium's adaptive response, and the vortex stability equation that follows from it. It has a rest density $\rho_0$,
a compressibility $\kappa_\text{vac}$, and a viscosity $\eta$. Its dynamics are
governed by the compressible Navier-Stokes equation:
\begin{equation}
  \rho\left(\frac{\partial \bm{v}}{\partial t}
  + (\bm{v}\cdot\nabla)\bm{v}\right)
  = -\nabla P + \eta\,\nabla\times(\nabla\times\bm{v}).
  \label{eq:NS}
\end{equation}
The left-hand side is the inertia of the medium --- density times acceleration.
The first term on the right is the pressure gradient that drives flow from
high-pressure to low-pressure regions. The second term is the viscous resistance
to rotational deformation. Nothing exotic here: this is the same equation that
governs water, air, and every other real fluid.

\subsection{The speed of light as a wave speed}

For small disturbances from the rest state, Eq.~\eqref{eq:NS} linearises.
Setting $\bm{v} = \bm{v}_0 + \delta\bm{v}$ and keeping only first-order terms
gives a wave equation for the displacement potential $\varphi$:
\begin{equation}
  \frac{\partial^2 \varphi}{\partial t^2}
  = \frac{\kappa_\text{vac}}{\rho_0}\,\nabla^2\varphi.
  \label{eq:wave}
\end{equation}
This is a standard compression wave. Its speed is $\sqrt{\kappa_\text{vac}/\rho_0}$.
CDFT identifies this with the speed of light:
\begin{equation}
  c = \sqrt{\frac{\kappa_\text{vac}}{\rho_0}}.
  \label{eq:c}
\end{equation}
This is a modelling postulate. It does not establish that electromagnetic waves
are compression waves in a material vacuum, and it needs an independent
discriminant beyond standard electrodynamics.

\subsection{Transport flux and the adaptive capacity constraint}

Define the transport flux density as $\bm{J} = \rho\bm{v}$ --- the amount of
medium moving through a surface per unit time per unit area. In any finite medium,
this cannot be infinite. There is a maximum flux $J_\text{crit}$ that the medium
can sustain before it saturates.

In the extended CDFT paradigm, this saturation is governed by the adaptive operating
ratio $\Psi_s$, the explicit form of the Mujjabi Adaptive Operating Ratio:
\begin{equation}
  \Psi_s = \left(\frac{|\bm{J}|}{J_\text{crit}}\right) \cdot S \cdot M_s \leq 1.
  \label{eq:Psi_s}
\end{equation}
Here, $S$ represents the surface responsiveness and $M_s$ represents the accumulated
surface memory. In the fundamental physical vacuum, which exists as a perfectly stable
equilibrium attractor, these adaptive variables converge to unity ($S \cdot M_s \to 1.0$).
Therefore, the operating ratio simplifies to the classical constraint:
\begin{equation}
  |\bm{J}| \leq J_\text{crit}.
  \label{eq:Jcrit}
\end{equation}
To see why this limit must exist, consider a region where flux is forced upward
by an external source. As $J \to J_\text{crit}$, the local pressure gradient
required to maintain flow diverges. The medium cannot comply. Instead of flowing
through, the flux is redirected \emph{around} the region. This redirection takes
the form of a vortex ring: the medium spirals around the overloaded core,
maintaining global transport while locally the flux is regulated. The vortex ring
is not inserted by hand --- it is the medium's natural response to approaching
its transport limit.


\section{Vortex Ring Geometry and Energy}
\label{sec:geometry}

\subsection{Two independent length scales}

A vortex ring in the transport medium is characterised by two lengths.

The first is the ring radius $R$, the distance from the axis of symmetry to the
centre of the vortex core. In CDFT, vortex circulation is quantised:
$\Gamma = \hbar/m$, where $m$ is the particle mass. The circulation also satisfies
the transport constraint: the tangential velocity at the ring surface equals $c$,
so $\Gamma = 2\pi R \cdot c$. Combining these:
\begin{equation}
  R = \frac{\Gamma}{2\pi c} = \frac{\hbar}{mc}.
  \label{eq:R}
\end{equation}
This is the reduced Compton wavelength. It arises here purely from quantum
circulation quantisation and the transport velocity constraint.

The second length scale is the core radius $a$, the size of the vortex core
itself. This is set by the electromagnetic self-energy condition: the electrostatic
energy stored in the field of a charge $e$ concentrated within radius $a$ must
equal the rest energy $mc^2$:
\begin{equation}
  \frac{e^2}{a} = mc^2
  \quad\Longrightarrow\quad
  a = \frac{e^2}{mc^2}.
  \label{eq:a}
\end{equation}
This is the classical electron radius~\cite{Jackson1999}. It arises from
electromagnetism and has nothing to do with quantum mechanics.

The key point is that $R$ and $a$ come from entirely different sectors of
physics --- quantum circulation on one side, electromagnetic self-energy on the
other. They are not connected to each other by construction. Yet their ratio will
turn out to be $1/\alpha$.

\subsection{The aspect ratio}

The dimensionless shape of the vortex ring is captured by the aspect ratio:
\begin{equation}
  \chi = \frac{R}{a}.
  \label{eq:chi}
\end{equation}
A large $\chi$ means a fat ring with a tiny core --- like a smoke ring. A small
$\chi$ means a tightly coiled ring where the core is almost as large as the ring
itself. The stability of the vortex, and its energy, depend entirely on $\chi$.

\subsection{The vortex ring energy}

The kinetic energy of a vortex ring moving through a fluid was derived by Kelvin
and systematised by Lamb~\cite{Lamb1932}. The result comes from integrating the
kinetic energy $\frac{1}{2}\rho v^2$ over all space, using the known velocity
field induced by a vortex ring of circulation $\Gamma$, radius $R$, and core
radius $a$. The integral evaluates to:
\begin{equation}
  E_\text{kin} = \rho_0 \Gamma^2 R
  \left[\ln\!\left(\frac{8R}{a}\right) - \beta\right],
  \label{eq:Ekin}
\end{equation}
where $\beta$ is a dimensionless constant that depends on how vorticity is
distributed inside the core. For a solid-body-rotation core --- vorticity
uniform across the cross-section, which Lamb~\cite{Lamb1932} showed is the
natural equilibrium distribution --- the integral gives:
\begin{equation}
  \beta = \tfrac{7}{4} = 1.75.
  \label{eq:beta}
\end{equation}
Substituting $R = a\chi$ and normalising by $2\pi\rho_0 a^3 c^2$, the
circulation energy becomes:
\begin{equation}
  E_\text{circ}(\chi) = \chi\left[\ln(8\chi) - \beta\right].
  \label{eq:Ecirc}
\end{equation}
This term grows with $\chi$: bigger rings have more kinetic energy. Left alone,
the vortex would expand indefinitely.

\subsection{The transport back-pressure term}

CDFT adds a second contribution to the energy that Kelvin's formula does not
contain. When the core radius $a$ is small (large $\chi$), the curvature at the
core boundary is high. High curvature forces the local transport flux toward
$J_\text{crit}$. The medium resists this: it generates a back-pressure that
opposes further tightening of the core. This back-pressure contributes an energy
cost that grows as the ring shrinks:
\begin{equation}
  E_\text{back}(\chi) = \frac{\kappa}{\chi},
  \label{eq:Eback}
\end{equation}
where $\kappa$ is a dimensionless transport stiffness parameter set by the vacuum
medium properties. This term diverges as $\chi \to 0$: the medium strongly resists
a vortex ring that is trying to collapse to a point.

\subsection{Total energy and the competition}

The full normalised vortex energy is the sum of these two contributions:
\begin{equation}
  \boxed{E_\text{norm}(\chi)
  = \chi\!\left[\ln(8\chi) - \beta\right] + \frac{\kappa}{\chi}.}
  \label{eq:Enorm}
\end{equation}

\begin{figure}[h]
  \centering
  \includegraphics[width=0.8\textwidth]{figures/fig1a_energy_landscape.pdf}
  \caption{The selected energy functional after $\kappa$ has been calibrated to place its minimum at $\chi^*=137.036$. This conditional minimum does not establish a physical vortex or fine-structure mechanism.}
  \label{fig:energy_landscape}
\end{figure}

The first term wants the ring to be large. The second term wants the ring to be
small. They compete. Somewhere between the two extremes, there is a minimum ---
a stable minimum of the chosen energy functional. Figure~\ref{fig:energy_landscape}
illustrates this conditional mathematical property; it does not identify a
particle or an observed vacuum state.


\section{The Stability Equation}

\subsection{Variational derivation}

The equilibrium condition follows from the variational principle. We minimise the
total energy subject to the transport constraint $J = J_\text{crit}$, enforcing
it with a Lagrange multiplier $\lambda$:
\begin{equation}
  \delta\!\left[E + \lambda(J - J_\text{crit})\right] = 0.
  \label{eq:variational}
\end{equation}
The Lagrange multiplier $\lambda$ has a clear physical meaning: it is the pressure
the medium exerts to prevent the constraint from being violated. At equilibrium,
$\lambda$ adjusts itself so that the vortex ring sits exactly at $J = J_\text{crit}$
without overshooting.

Differentiating $E_\text{norm}(\chi)$ with respect to $\chi$:
\begin{align}
  \frac{dE_\text{norm}}{d\chi}
  &= \frac{d}{d\chi}\left\{\chi\!\left[\ln(8\chi)-\beta\right]
     + \frac{\kappa}{\chi}\right\} \nonumber\\
  &= \left[\ln(8\chi) - \beta\right]
     + \chi\cdot\frac{1}{\chi}
     - \frac{\kappa}{\chi^2} \nonumber\\
  &= \ln(8\chi) - \beta + 1 - \frac{\kappa}{\chi^2}.
  \label{eq:dEdchi}
\end{align}
Setting this to zero gives the equilibrium condition:
\begin{equation}
  \boxed{\ln(8\chi) - \beta + 1 = \frac{\kappa}{\chi^2}.}
  \label{eq:equil}
\end{equation}

\begin{figure}[h]
  \centering
  \includegraphics[width=0.8\textwidth]{figures/fig1b_equilibrium.pdf}
  \caption{Equilibrium condition $dE/d\chi = 0$ after the stiffness has been fitted to the target $\chi=1/\alpha$.}
  \label{fig:equilibrium}
\end{figure}

\subsection{Confirmation that it is a minimum}

An equilibrium could be a minimum (stable) or a maximum (unstable). To confirm
stability, we check the second derivative at $\chi = 137$:
\begin{equation}
  \frac{d^2 E_\text{norm}}{d\chi^2}
  = \frac{1}{\chi} + \frac{2\kappa}{\chi^3}.
  \label{eq:d2E}
\end{equation}
Both terms are strictly positive for all $\chi > 0$, so the selected energy
functional has a stable minimum. This calculus result is conditional on the
postulated back-pressure term and does not determine its physical parameter.

\subsection{The required transport stiffness}

Rearranging Eq.~\eqref{eq:equil} gives the transport stiffness $\kappa$ that
places the energy minimum at any desired $\chi$:
\begin{equation}
  \kappa(\chi) = \chi^2\left[\ln(8\chi) - \beta + 1\right].
  \label{eq:kappa_for_chi}
\end{equation}
At the pre-specified target $\chi=1/\alpha=137.036$, this evaluates to
$\kappa=117{,}362$. It is the fitted dimensionless parameter required by this
functional, not a determined vacuum transport stiffness. No independent route
from $\rho_0$ and $c$ to this value has been established. Figure~\ref{fig:sensitivity}
only shows local sensitivity of the calibrated functional.

\begin{figure}[h]
  \centering
  \includegraphics[width=0.8\textwidth]{figures/fig1c_sensitivity.pdf}
  \caption{Sensitivity of the calibrated energy minimum to variations in the fitted parameter $\kappa$. It is a local property of the selected functional, not a prediction of fine-structure fluctuations.}
  \label{fig:sensitivity}
\end{figure}


\section{Definitional Relation of $\chi$ and $\alpha$}

The following calculation makes explicit why the advertised geometric agreement
is not independent evidence.

From quantum circulation quantisation and $v(R) = c$, we obtained:
\begin{equation*}
  R = \frac{\hbar}{mc} = 3.8616 \times 10^{-13}\ \text{m.}
\end{equation*}
From electromagnetic self-energy $e^2/a = mc^2$, we obtained:
\begin{equation*}
  a = \frac{e^2}{mc^2} = 2.8179 \times 10^{-15}\ \text{m.}
\end{equation*}
Their ratio is:
\begin{equation}
  \chi = \frac{R}{a}
  = \frac{\hbar/mc}{e^2/mc^2}
  = \frac{\hbar c}{e^2}
  = \frac{1}{\alpha}.
  \label{eq:chi_is_alpha}
\end{equation}
Every cancellation is shown. Nothing is hidden. The $mc$ factors cancel, leaving
the combination $\hbar c / e^2$, which is exactly the inverse of the fine-structure
constant by definition.

This is an algebraic identity. The standard reduced Compton wavelength and
classical electron radius contain the same constants that define $\alpha$.
Dividing them therefore recovers $1/\alpha$ by construction; it neither tests
the vortex ansatz nor supplies an additional physical constraint.


\section{Numerical Results}
\label{sec:results}

All computations were performed using the paper-local supplementary script and
notebook listed in the Data and Code Availability section. These public
artifacts implement only the equations stated in this manuscript, cross-check
roots and integrals with SciPy, verify symbolic derivatives with SymPy, and
export reproducible CSV outputs with Pandas.

\subsection{Stability equation}

Using $\beta = 1.75$ and the target $\chi_\text{target} = 1/\alpha = 137.035999177$
(CODATA 2022~\cite{CODATA2022}), Eq.~\eqref{eq:kappa_for_chi} gives:
\begin{equation*}
  \kappa = 117{,}362.00.
\end{equation*}
We then solve Eq.~\eqref{eq:equil} numerically using bisection over
$\chi \in (1.01,\, 10^6)$, with $\kappa$ fixed to this value. The recovered
equilibrium is:
\begin{equation*}
  \chi_\text{recovered} = 137.035999177,
  \quad
  \left|\frac{\chi_\text{rec} - \chi_\text{target}}{\chi_\text{target}}\right|
  = 1.19 \times 10^{-12}.
\end{equation*}
The fitted functional has a stable solution at the target $\chi=1/\alpha$. The
reported difference is numerical precision --- the bisection algorithm's
convergence tolerance --- rather than a prediction error.

\begin{table}[h]
  \caption{Calibration input and root-solver output for the same equilibrium equation.}
  \label{tab:stability}
  \begin{tabular}{lc}
    \toprule
    Quantity & Value \\
    \midrule
    $\chi_\text{target} = 1/\alpha$ & $137.035999177$ \\
    $\beta$ (Lamb solid-rotation) & $1.75$ \\
    $\kappa$ fitted to target & $117{,}362.00$ \\
    $\chi_\text{recovered}$ (bisection) & $137.035999177$ \\
    $\alpha_\text{recovered} = 1/\chi$ & $0.007297352566$ \\
    $\alpha_\text{measured}$ (CODATA 2022) & $0.007297352566$ \\
    Root-solver tolerance & $1.19\times10^{-12}$ \\
    \bottomrule
  \end{tabular}
\end{table}

\subsection{Energy balance}

Table~\ref{tab:energy} shows the two energy contributions at the equilibrium point.

\begin{table}[h]
  \caption{Energy balance at $\chi = 137.036$, in units of $2\pi\rho_0 a^3 c^2$.}
  \label{tab:energy}
  \begin{tabular}{lc}
    \toprule
    Quantity & Value \\
    \midrule
    $E_\text{circulation}$ & $719.40$ \\
    $E_\text{back-pressure}$ & $856.43$ \\
    Ratio $E_\text{back}/E_\text{circ}$ & $1.19$ \\
    \bottomrule
  \end{tabular}
\end{table}

\begin{figure}[h]
  \centering
  \includegraphics[width=0.8\textwidth]{figures/fig1d_energy_components.pdf}
  \caption{The individual energy components ($E_\text{circ}$ and $E_\text{back}$) plotted against the total energy. Their precise intersection behaviour generates the shallow potential well.}
  \label{fig:energy_components}
\end{figure}

The ratio is 1.19 within the calibrated functional. This numerical balance says
nothing about whether the fitted parameter is natural or independently fixed.

\subsection{Definitional length-ratio identity}

Table~\ref{tab:selfconsist} evaluates the definitional identity in
Eq.~\eqref{eq:chi_is_alpha}, using CODATA 2022 values.

\begin{table}[h]
  \caption{Definitional identity $R/a=1/\alpha$ evaluated from standard length scales.}
  \label{tab:selfconsist}
  \begin{tabular}{lc}
    \toprule
    Quantity & Value \\
    \midrule
    $R$ (reduced Compton wavelength) & $3.861593\times10^{-13}$ m \\
    $a$ (classical electron radius) & $2.817940\times10^{-15}$ m \\
    $\chi_\text{geom} = R/a$ & $137.035999177$ \\
    $\alpha_\text{geom} = a/R$ & $0.007297352564$ \\
    $\alpha_\text{measured}$ & $0.007297352564$ \\
    Rounding difference & $2.26\times10^{-12}$ \\
    \bottomrule
  \end{tabular}
\end{table}

The agreement follows from the definitions and is limited only by the rounding
of the input constants; it is not an independent model check.


\section{Discussion}

\subsection{Interpretation boundary}

The fine-structure constant remains an empirical QED coupling. This manuscript
shows only that a chosen energy functional can be fitted to have a minimum at
the known value and that standard length-scale definitions algebraically recover
the same value. Those facts do not identify $\alpha$ as a vortex aspect ratio or
provide a vacuum mechanism.

\subsection{On the naturalness of $\kappa$}

Equation~\eqref{eq:kappa_for_chi} makes the calibration explicit. Its scaling
with the target does not determine $\kappa$ independently. A previously proposed
dimensional form was
\begin{equation}
  \kappa = \frac{\kappa_\text{vac}}{\rho_0 c^2 a},
  \label{eq:kappa_dim}
\end{equation}
but substituting Eq.~\eqref{eq:c} gives $\kappa=1/a$, not $1/(a\rho_0)$.
The unresolved dimensional mismatch means this expression cannot supply a
physical closure. Paper~III did not establish an independent derivation.

\subsection{Memory, hysteresis, and the limits of the equilibrium calculation}

The adaptive factors $S$ and $M_s$ are not used as independent fitting knobs in
the numerical recovery above. The calculation assumes its hypothetical vacuum
medium is in an equilibrium limit, so $S\cdot M_s\rightarrow1$ and
Eq.~\eqref{eq:Psi_s} reduces to the capacity condition. The purpose of retaining
$S$ and $M_s$ in the notation is to mark the non-equilibrium extension that later
papers must test rather than to improve the fit in this paper.

A minimal way to state the hypothesis is to let vacuum memory be a local
relaxing state variable:
\begin{equation}
  M_s(\bm{x},t)
  =
  1+\mu\int_{-\infty}^{t}
  \max\!\left(0,\frac{|\bm{J}(\bm{x},t')|}{J_\text{crit}}-1\right)
  e^{-(t-t')/\tau_M}\,dt'.
  \label{eq:memory_kernel}
\end{equation}
Here $\tau_M$ is a candidate relaxation time and $\mu$ measures how strongly
over-capacity flux leaves residual constraint. In ordinary low-flux settings,
$|\bm{J}|\ll J_\text{crit}$ and the integral vanishes, so $M_s\approx1$ and the
vacuum appears memoryless. In extreme flux settings, a nonzero $M_s-1$ would mean
that the medium has not fully reset. This is a physical hypothesis, not a result
already established by the present calculation.

This distinction matters for quantum decoherence. Within standard quantum
mechanics, decoherence is described by entanglement with environmental degrees of
freedom. CDFT can only add value if the proposed memory variable predicts a
measurable residual phase, relaxation, or hysteresis effect beyond the standard
description. Paper~XII states this as an experimental target rather than as an
assumed fact.

This is also the intended meaning of any observer-effect language in this
release. An observer is not a disembodied human mind changing the vacuum by
attention. Operationally, an observation is a physical coupling between the
system and an apparatus, environment, detector, or record. If such coupling
changes the state, the mechanism must be specified as an interaction channel.
CDFT therefore treats ``observation'' as a boundary-and-measurement problem, not
as a consciousness claim.

\subsection{Operational dictionary relative to standard quantum language}

The source notes for this release also proposed a comparison between the
transport-medium vocabulary and standard quantum language. The professional way
to use that comparison is not as a replacement of established theory, but as an
operational dictionary: it states what CDFT would have to mean in measurable
terms if the adaptive-vacuum interpretation is physical.

\begin{center}
\scriptsize
\begin{tabular}{p{0.16\textwidth}p{0.21\textwidth}p{0.31\textwidth}p{0.14\textwidth}}
\toprule
\textbf{Feature} & \textbf{Standard reading} & \textbf{CDFT reading} & \textbf{Status}\\
\midrule
Vacuum & Background state with quantum fields & Effective transport medium with finite capacity $C$ and response variables $S,M_s$ & Hypothesis\\
Fine-structure constant & Empirical QED coupling & Candidate vortex-aspect-ratio interpretation $\chi=R/a$ & Calibrated toy model\\
Wave--particle duality & Complementary measurement descriptions & Smooth flux transport below threshold, vortex regulation near $\Psi_s\approx1$ & Interpretation\\
Decoherence & Environmental entanglement and phase suppression & Possible residual vacuum-memory lag in addition to standard decoherence & Test target\\
Planck constant & Fundamental quantum of action & Candidate product of a critical flux threshold and vortex period, $h=\Phi_c\tau_v$ & Open derivation\\
Charge and mass & Empirical quantum numbers and inertial scale & Candidate boundary discontinuity and inertial resistance of a constrained vortex state & Open mapping\\
Gravity bridge & Spacetime curvature / Newtonian attraction & Possible secondary pressure-gradient limit of constrained transport & Speculative bridge\\
\bottomrule
\end{tabular}
\end{center}

This table is deliberately conservative. It does not claim that CDFT has
replaced QED or general relativity. It identifies the places where the
transport-medium model must either produce new predictive value or fail.

\subsection{Core-profile robustness program}

The twelve-decimal recovery is an internal check of the stated equilibrium
equation. It is not yet a proof that the same fixed point survives every possible
vortex-core model. A professional robustness program must vary the internal
vorticity profile and ask whether the same aspect ratio is recovered without
retuning the transport stiffness:
\begin{center}
\small
\begin{tabular}{p{0.22\textwidth}p{0.28\textwidth}p{0.32\textwidth}}
\toprule
\textbf{Core model} & \textbf{Why it matters} & \textbf{Required test}\\
\midrule
Rankine / solid core & Baseline Lamb-style compact vortex model & Verify that the existing $\beta=7/4$ calculation is stable under resolution and derivative checks.\\
Gaussian / diffuse core & Mimics smeared or soft-boundary flux density & Recompute $\beta$ and the energy minimum with a width convention fixed before comparison.\\
Hollow or excluded core & Tests the extreme capacity-limited boundary case & Determine whether the back-pressure term changes functional form.\\
Finite-thickness vortex rings & Tests departure from the thin-ring approximation & Compare against finite-core ring families before claiming topological invariance.\\
\bottomrule
\end{tabular}
\end{center}
Until those tests are implemented in the public supplementary code, the correct
claim is conditional: the present model has a stable calibrated fixed point at
$\chi=1/\alpha$, and core-profile invariance is an open robustness target.

The current supplementary script now writes an explicit guardrail table and a
matching figure in the Paper I output directory.
Those outputs do not prove invariance. They show how alternative core
assumptions move the equilibrium when the Rankine calibration is held fixed,
which is the correct starting point for the robustness program.

\subsection{Connection to the lepton masses}

This paper makes no particle-family assignment. Paper~II studies a selected
$\mathbb{Z}_3$ parametrisation and its relation to the Koide mass formula,
\begin{equation*}
  Q = \frac{m_e + m_\mu + m_\tau}
           {(\sqrt{m_e} + \sqrt{m_\mu} + \sqrt{m_\tau})^2}
  = \frac{2}{3},
\end{equation*}
which holds experimentally to six significant figures. That conditional
parametrisation is not evidence that the electron, muon, and tau are vortex
orientations, and it supplies no physical particle assignment.


\section{Necessary Conditions for the Model to Hold}

The alpha recovery in this paper is conditional. The following conditions must
hold for the CDFT interpretation to be more than a calibrated numerical
coincidence.

\begin{enumerate}
\item \textbf{Topological resilience.}
The aspect ratio $\chi=R/a$ must be a stable eigenmode, or fixed point, of the
adaptive operating ratio $\Psi_s$. In practical terms, small perturbations of
the ring geometry must relax back toward $\chi^*$ rather than drift to a
continuum of nearby values. If perturbations select many unrelated aspect
ratios with equal stability, the fine-structure interpretation fails.

\item \textbf{Flux--memory coupling.}
The responsiveness $S$ and memory $M_s$ must be coupled so that the equilibrium
vacuum satisfies $S M_s\to1$ in the low-flux limit while still allowing a
finite, measurable residual after overload. The simplest admissible behavior is
an inverse correlation: if $M_s$ rises after a high-flux event, the effective
response must relax so that the attractor is restored instead of producing
unbounded amplification.

\item \textbf{Non-radiative equilibrium.}
The geometry $\chi^*\approx137.035999177$ must represent a minimum of internal
dissipation for the regulator vortex. In the language of this model, the
equilibrium is a superfluid-like or zero-friction transport state; outside the
minimum, excess curvature or excess ring size should introduce loss, radiation,
or instability. This is a condition to be tested, not a claim already
established by the twelve-decimal calculation.

\item \textbf{Independent stiffness closure.}
The stiffness $\kappa$ must be derived from the vacuum-medium properties without
feeding in $\alpha$, $e$, or $\hbar$ as hidden calibration targets. Until that is
done, the Paper~I result remains a high-precision consistency closure rather
than a parameter-free derivation.

\item \textbf{External discriminant.}
At least one measurable prediction must survive comparison with established
physics. Candidate tests include vacuum hysteresis, extreme-field spectral
residuals, decoherence-memory residuals, rotating-source pressure-gradient
residuals, and transport-threshold tests in collective matter.
\end{enumerate}

These conditions are intentionally strict. They are the difference between a
suggestive geometric reconstruction of $\alpha$ and a physical theory that can
be rejected by experiment.


\section{Open Problems}

We state the following explicitly, because we believe clear honesty about what
remains to be done is more valuable than overclaiming:

\begin{enumerate}

\item \textbf{Find an independent physical closure for $\kappa$.}
The transport stiffness $\kappa = 117{,}362$ was computed from the condition
that the energy minimum sits at $\chi = 137$. It has not yet been derived
independently from the vacuum medium properties $\rho_0$ and $c$. The proposed
Paper~III route did not establish that derivation. Until an independently
dimensionally consistent closure exists, this model cannot predict $\alpha$.

\item \textbf{Specify an independently testable route for $\hbar$.}
We used $\Gamma = \hbar/m$ to derive $R$. A complete derivation would obtain
$\hbar$ from the vacuum quantisation rule $\Gamma = n\Gamma_0$ without using
$\hbar$ as an input. A concrete target suggested by the source notes is
\begin{equation}
  h=\Phi_c\,\tau_v,
\end{equation}
where $\Phi_c$ is a critical vacuum-flux threshold and $\tau_v$ is the
characteristic period of the vortex attractor. At present this is a target
identity, not a completed derivation or a result of Paper~III.

\item \textbf{Specify an independently testable route for $e$.}
We used $e^2/a = mc^2$ to derive $a$. A complete derivation would obtain $e$
from the condition that electromagnetic flux is quantised in the medium. In the
same language, charge would be the measurable discontinuity induced at the
adaptive flux boundary rather than an unexplained label attached to the vortex.

\item \textbf{Test core-profile robustness without retuning $\kappa$.}
The current public calculation uses the stated Lamb-style core model. A stronger
derivation must repeat the stability calculation for hard, diffuse, hollow, and
finite-thickness profiles with conventions fixed in advance.

\item \textbf{Measure or bound the vacuum-memory relaxation time $\tau_M$.}
The memory kernel in Eq.~\eqref{eq:memory_kernel} is a hypothesis. It becomes
physics only if an experiment or a deeper theory can identify, bound, or rule out
the relaxation time and coupling strength.

\item \textbf{Find an external discriminant.}
The recovery of $1/\alpha$ is a postdiction of an already known constant unless
the framework also predicts a measurable effect not already accounted for by QED,
general relativity, or standard condensed-matter and plasma models.

\end{enumerate}

Until an independently derived, dimensionally consistent parameterisation and a
new discriminating observation exist, the experimental value of $\alpha$ remains
only the calibration target used to inspect this model's internal algebra.


\section{Conclusion}

This paper documents a calibrated toy energy functional. Choosing $\kappa$ from
the known target $\chi=1/\alpha$ and solving the same equilibrium equation
returns that target to the numerical tolerance of the solver. The ordinary
relation $R/a=\hbar c/e^2=1/\alpha$ is a definitional identity. Neither is an
independent prediction or evidence that the physical vacuum is a constrained
transport medium.

The conditional second-derivative calculation remains valid for the specified
functional. The physical mapping, the proposed back-pressure term, the
dimensions and independent origin of $\kappa$, and every vacuum-memory extension
remain open hypotheses. This paper must not be cited as a derivation of the
fine-structure constant or as support for downstream universal claims.


\section*{Supplementary Material}

This manuscript is accompanied by release-archive supplementary material in
the canonical Part I tree. The relevant paper-local Python scripts, numerical checks, generated figure outputs,
tabulated data, figure assets, and environment notes are maintained under
\texttt{notebooks/}, \texttt{outputs/}, \texttt{figures/}, and
\texttt{REPRODUCIBILITY.md} in \texttt{Part\_I\_Fundamental\_Physics/}.


\section*{Data and Code Availability}

The release-archive supplementary material is included in the archived Part I
release on Zenodo at \url{https://doi.org/10.5281/zenodo.20250820}. The current version DOI is assigned by Zenodo after each GitHub release. Zenodo is the
permanent release archive, and the public source archive is available on GitHub at
\url{https://github.com/bampita-bico/CDFD-Part-I-Release}. The broader public CDFD
Runtime is separate reusable software infrastructure; the numerical values reported
here are reproduced from this paper's Supplementary Material.


\section*{Acknowledgments}
The author used AI-assisted tools for formatting checks and code-verification
support. The scientific claims, interpretations, and final manuscript decisions
are the author's responsibility.

The author is indebted to the foundational work of
Lord Kelvin~\cite{Thomson1867} (vortex atom hypothesis),
H.~Lamb~\cite{Lamb1932} (vortex ring energy),
P.G.~Saffman~\cite{Saffman1992} (vortex dynamics),
H.~von Helmholtz~\cite{Helmholtz1858} (vortex theorems),
and P.A.M.~Dirac~\cite{Dirac1951} (vacuum as physical medium)
--- whose ideas this paper attempts to build upon.

\section*{Declarations}

\subsection*{Funding}
This research received no external institutional funding.

\subsection*{Author Contributions}
Steve Bico Mujjabi: conceptualization, methodology, formal analysis, software,
validation, visualization, writing--original draft, and writing--review and
editing.

\subsection*{Conflicts of Interest}
The author declares no conflicts of interest.

\section*{References}
\bibliographystyle{plainnat}
\bibliography{references}

\end{document}
